% !TeX root = ../papers/nc-p3.tex
\section{四元群を仮定しない十四次元族と二つの局所変形}
\label{sec:unrestricted-doublet}
$r$を節\ref{sec:klein-generic-isomorphisms}の反対階数10の有理点、
$t$を節\ref{sec:klein-elliptic-completion}の反対階数16の有理正則点とする。
$r$の正則性にも四次中心の証明書がある。ここでは両点で全ての文字次数の
変形を許す。自己同型は四位置の射影基底変換で測る。

\begin{theorem}\label{thm:unrestricted-doublet}
$r$を通る片側対称な窓には、添字保持同型を除いて十四次元の既約族がある。
幾何学的一般点は反対階数10、最小周期4の標準Koszul左右AS正則窓で、
その射影自己同型群は自明である。片側対称性まで外した軌道に垂直な接空間は
$r,t$ともに十五次元で、両点で全ての二次障害が消える。
$r$における片側対称族の外への一方向は、反対切断の六次元の核から余核への同型を誘導する。
\end{theorem}

\begin{proof}
$W=(w_{ij,kl})$、$W^\tau=(w_{kj,il})$とすると、二切断条件はこれらの階数が
共に6となる条件である。各切断で六次の可逆な小行列を選び、その逆を逆順の
行・列位置に埋めた行列を$U$と書く。左核・右核の基底行列を$L,R$とすれば、
Schur補行列の一次条件は$D(T)=L^{\mathsf t}TR=0$である。
二方向$S,T$の二次障害は$D(SUT+TUS)$を一次条件の像で割った類になる。
二切断を同時に扱うとき、二番目にはその切断自身の核と逆行列を使う。
これは可逆小行列の逆の二次までの展開から従う。標数は零である。

$r$では$W^\tau=W$なので、交換の対称・交代部分に分ければ、二切断の一次条件は
それぞれ単独の$D$になる。符号の固有空間をさらに四元群の文字で分け、
四つの位置の無限小基底変換を除いた結果を示す。各組は文字次数$0,1,2,3$の順である。
\[
\begin{array}{c|c|c}
 &\text{アフィン軌道方向}&\text{軌道を除いた接方向}\\\hline
r,\ \tau=+1 &(10,12,12,12)&(4,2,4,4)\\
r,\ \tau=-1 &(3,4,4,4)&(1,0,0,0)\\
t,\ \text{全方向}&(13,16,16,16)&(5,2,4,4)
\end{array}
\]
計算には$ i+j+k+l$が指定文字となるテンソル単位を用いる（加法は四元群内）。
$r$では偶二位置を交換した和または差を取り、対角項は一度だけ使う。
各$\mathfrak{gl}_4$の行列単位を一位置に作用させれば軌道方向を得る。
\path{computations/verify_unrestricted_doublet_deformations.py}は全てを有理数上で
掃き出し、軌道方向が実際に核に入ることと、この表の各階数を検算する。
対称部分の160変数から百個のSchur条件への微分は全射である。
よって固定した偶交換を持つ階数6軌跡は$r$で滑らかな60次元である。
反対像が対称二乗全体なので、二つの片側対称表示の間の偶二位置の基底変換は
射影的に一致する。実際、任意の平方の像が再び対称になるには二写像が各直線上で比例する。
従ってこの軌跡上の一般の軌道次元は46（アフィン倍率込み）で、
有理商の次元は$60-46=14$である。軌道の有限な射影安定群は以下からも従う。

三次、四次の縮約が作る有限Frobenius双対に対し、節\ref{sec:reversible-component}の
内部次数別の特殊化をこの既約族と正則点$r$に適用すると、幾何学的一般点で全次数の
正則性を得る。反対階数10は片側対称部分で開く。異なる二奇配置が互いに射影非同値で
あることは$r$での有限被約な二十点配置の比較により一般点に保たれるため、最小周期は4である。

同じ有限エタールな点の番号付けを$r$の近くで施す。一般点で窓を保つある置換が存在すれば、
各位置でその二十点を運ぶ行列の存在条件は射影的に閉で、$r$に特殊化する。
三点直線・五点平面がないため極限行列は可逆である
（節\ref{sec:klein-generic-isomorphisms}）。特殊点で各位置に五点の射影枠を選ぶと
候補行列は枠から一意かつ正則に復元され、窓を保つ等式は既約な近傍に延びる。
特殊点における窓の射影安定群は対角四元群である。実際、各配置の安定群は四元群で、
辺が点集合間のグラフであるため一位置の作用が全ての位置の作用を決める。
特殊値を$g$として、変動する行列によるテンソルの固定等式を微分すると、
任意の接ベクトル$v$について$gv-v$は軌道接空間に入る。番号付けはエタールなので
全接方向でこれを要求できる。しかし上表の対称行には全ての非自明文字が現れ、
$g$が商接空間を点ごとに固定すれば$g=1$である。有限エタールな番号付けの
異なる置換は特殊点でも異なるので、一般点の置換も恒等に限られる。

$r$の二次条件では対称同士と交代同士の積が対称側に入るので、上の全射で解ける。
交代方向の代表一つと十四の対称代表の混合積は、各文字の交代微分の像に入る。
$t$では十五代表の重複を許す全120組を二切断で同時に調べる。
どちらも上記コードで障害列を付け加えた階数増分が全て零となる。
基底変換自体は持ち上がるので軌道方向を除く計算で十分である。
最後に$r$で反対切断の左右核で唯一の交代代表を挟むと階数6を得る。
従ってその接線を持つ弧があれば反対行列式の最低次はちょうど6である。
\end{proof}

\paragraph{階数条件を保つ無限小変形の平坦性。}
二つの切断の階数が局所自由な意味で6である局所Artin環上の持ち上げで、
特殊ファイバーが標準AS窓なら、その四窓代数は各内部次数で自由となる。
実際、巡回したテンソル、末尾一文字と二文字を縮約した像、矢空間、基底環を
順に使い、各始点に階数$1,4,6,4,1$の係数を持つ線型複体を作る。
これらの像は単元小行列で直和因子となり、連続二縮約の積は二次関係に入るので
微分の合成は零である。内部次数$n$で、正のホモロジー位置の項に現れる
代数の次数は$n$未満である。帰納的にそれらは自由で、特殊ファイバーは
Koszul分解なので正の位置で完全である。左端から単元小行列を順に分離すると
基底環上でも正の位置で完全で、矢から掛け合わせた次数$n$商は自由になる。
より具体的に末尾の商は関係を挿入した
$R_i\otimes A_{i+2,i+n}\to V_i\otimes A_{i+1,i+n}$の余核であり、
二次表示の定義通り$A_{i,i+n}$である（$n=1$では関係項は零）。
左側も逆順で同じ証明となり、traceの相補次数対合は特殊化で完全なので完全である。
この議論は二次までの持ち上げにも適用される。十五の一次方向が全て同時に
任意次数まで持ち上がるかは、二次障害の消滅だけでは決まらない。
